\section{Results}
\label{sec:results}

\subsection{H1: the ground truth decides the ranking}
\label{sec:res-h1}
Table~\ref{tab:dualgt} scores every model under every ground truth. With the U-Net, the \GTb-trained model reaches a pooled AUPRC of 0.923 under \GTb{} but 0.374 under \GTp{} and 0.257 under \GTg; on the official test set the same predictions give 0.950 under \GTb{} and 0.427 under \GTp{} ($n{=}89$ pairs). Each ground truth ranks the model trained on it first, and the \GTb{} ranking is the reverse of the \GTg{} ranking. DeepLabV3+ reproduces the pattern with nearly identical numbers (0.939 $\rightarrow$ 0.382), so the effect is a property of the supervision, not of the U-Net decoder. Under its own metric the \GTb{} model is the best localizer in the study; under a pixel-change metric it is worse than a model that never saw a box. Most of the 0.92 is therefore earned by labelling unchanged pixels inside the rectangle.

\begin{table}[t]
\caption{Pooled pixel AUPRC on validation (487 images) for each training mask (rows) scored against each ground truth (columns); mean$\pm$sd over three seeds where available. Bold: best per column within an architecture.}
\label{tab:dualgt}
\centering\footnotesize
\setlength{\tabcolsep}{4pt}
\begin{tabular}{@{}llccc@{}}
\toprule
Model & Train & \GTb & \GTp & \GTg \\
\midrule
\multirow{3}{*}{U-Net (3 seeds)} & \GTb & \textbf{\pmm{0.923}{.015}} & \pmm{0.374}{.008} & \pmm{0.257}{.007} \\
 & \GTp & \pmm{0.579}{.014} & \textbf{\pmm{0.610}{.012}} & \pmm{0.511}{.010} \\
 & \GTg & \pmm{0.351}{.015} & \pmm{0.379}{.012} & \textbf{\pmm{0.676}{.019}} \\
\midrule
\multirow{3}{*}{DeepLabV3+} & \GTb & \textbf{0.939} & 0.382 & 0.256 \\
 & \GTp & 0.537 & \textbf{0.583} & 0.432 \\
 & \GTg & 0.345 & 0.401 & \textbf{0.549} \\
\midrule
YOLO11s (3 seeds) & \GTb & \textbf{\pmm{0.965}{.003}} & \pmm{0.368}{.015} & \pmm{0.221}{.003} \\
\midrule
U-Net, marker-repaired & \GTb & 0.916 & 0.389 & 0.266 \\
 & \GTp$_{\mathrm{nm}}$ & 0.513 & 0.544 & 0.455 \\
 & \GTg$_{\mathrm{nm}}$ & 0.348 & 0.371 & 0.661 \\
\midrule
U-Net + box aux. & \GTp & 0.647 & 0.622 & 0.518 \\
 & \GTg & 0.370 & 0.391 & 0.636 \\
\midrule
U-Net, 20 epochs & \GTb & 0.926 & 0.373 & 0.259 \\
 & \GTp & 0.568 & 0.594 & 0.488 \\
 & \GTg & 0.398 & 0.390 & 0.645 \\
\bottomrule
\end{tabular}
\end{table}

\subsection{Shortcut and trace controls (H2, H4)}
\label{sec:res-controls}
The \GTb{} model predicts rectangles (median rectangularity 0.87 vs.\ 0.32 for \GTp{} and 0.05 for \GTg) and is indifferent to blurring the box interior (relative AUPRC change $+1.2\,\%$) while a 5-px blur of the box rim costs $-2.0\,\%$ (placebo ring 0.0\,\%; three seeds): it uses the field's extent more than its content, but the boundary seam is only a small cue (H2 partly supported). A ResNet-18 trained on 224-px crops that never touch the box separates forged from paired authentic images at chance (AUROC 0.516; 0.980 for box-centred crops through the same pipeline), so no global re-encoding trace is exploitable (H4 rejected). Benign paste-back and full-image resampling, which composite or resample without changing content, reduce image-level AUROC (99th-percentile score) by 1.6--3.0 points for the three-seed U-Nets and by at most 4.5 points for the other standard segmenters, but by 3.5--6.5 points for the marker-repaired models (one seed; the detector was not evaluated on these controls), and the \GTp{} model raises the fewest false alarms on authentic images (12\,\% vs.\ 26\,\% of images with $>$0.1\,\% positive pixels at 0.5; three seeds).

\subsection{The marker shortcut, and retraining without it}
\label{sec:res-marker}
Fig.~\ref{fig:probes} defines the input probes; Table~\ref{tab:marker} reports them on the 380 validation pairs whose forgery carries marker residue (three-seed means for the standard U-Nets, seed 42 for all other rows). For the standard box U-Net, copying only the residue pixels back from the authentic image lowers per-image AP from 0.988 to 0.766 ($\Delta{=}-0.222$ \ci{-0.256}{-0.190}), and pasting only the residue onto the authentic image makes the model cover 80\,\% of the box at $\tau^*$ (AP 0.905). DeepLabV3+ and the pixel-trained U-Net show the same dependence ($\Delta$ $-0.174$ and $-0.247$). On the official test set the box model's AP is 0.972 on the 627 images with residue and 0.690 on the 185 without (0.722 on the 102 residue-free Gemini images), so the gap is not only a generator gap.

Retraining on marker-repaired images removes the dependence without destroying the benchmark score. The repaired box model loses nothing when the residue is removed at test time ($\Delta{=}+0.019$ \ci{0.002}{0.036}); the residue alone now covers 42\,\% of the box instead of 80\,\% (11\,\% instead of 42\,\% for the pixel model); its test AUPRC is 0.939 (0.950 before), its AP on residue images falls from 0.972 to 0.920, while on residue-free Gemini images it is similar (0.756 vs.\ 0.722, overlapping CIs). The marker is thus a shortcut the models take when it is offered, not a prerequisite for high box scores. All repaired results use one seed.

\begin{table}[t]
\caption{Marker probes on 380 validation pairs with residue (standard U-Net rows: mean of three seeds; all other rows: seed 42). AP: per-image AP under \GTb{} with 95\,\% group-bootstrap CI for the repair contrast; Cover: fraction of the box above $\tau^*$; AUROC: forged-type input vs.\ authentic, maximum score.}
\label{tab:marker}
\centering\footnotesize
\setlength{\tabcolsep}{2.2pt}
\begin{tabular}{@{}lcccccc@{}}
\toprule
 & \multicolumn{2}{c}{forged} & repaired & $\Delta$ AP & \multicolumn{2}{c}{auth.+marker} \\
\cmidrule(lr){2-3}\cmidrule(lr){6-7}
Model & AP & Cover & AP & (repair) & Cover & AUROC \\
\midrule
U-Net \GTb & 0.988 & 0.94 & 0.766 & $-0.222$ & 0.80 & 0.903 \\
U-Net \GTp & 0.747 & 0.54 & 0.500 & $-0.247$ & 0.42 & 0.858 \\
DeepLabV3+ \GTb & 0.982 & 0.92 & 0.807 & $-0.174$ & 0.78 & 0.899 \\
U-Net+aux \GTp & 0.771 & 0.59 & 0.597 & $-0.174$ & 0.48 & 0.873 \\
\midrule
repaired \GTb & 0.939 & 0.89 & 0.958 & $+0.019$ & 0.42 & 0.695 \\
repaired \GTp$_{\mathrm{nm}}$ & 0.648 & 0.44 & 0.686 & $+0.038$ & 0.11 & 0.592 \\
\bottomrule
\end{tabular}
\end{table}

\begin{figure}[t]
\centering
\includegraphics[width=\columnwidth]{fig/fig_probe_conditions.png}
\caption{Input probes on one validation pair (dashed: official box). The forged field carries a 1-px green residue of the generation marker on its top edge; \emph{render-only} keeps the generator's re-rendering but restores the original digits; \emph{content-only} carries the changed digits into the authentic image without the re-rendering or the marker.}
\label{fig:probes}
\end{figure}


\subsection{Re-rendering versus content}
\label{sec:res-render}
Generators re-render the whole field, so inside the box many pixels change by less than the \GTp{} threshold. Separating the two components on all validation pairs (Table~\ref{tab:render}) shows that neither alone reproduces the response to a real forgery. For the standard box U-Net (three seeds) the mean probability inside the box falls from 0.86 (forged) to 0.24 (render-only) and 0.35 (content-only), and forged-vs-authentic AUROC from 0.971 to 0.637 and 0.733: the standard models respond most strongly to the conjunction of re-rendering, changed glyphs and marker. Removing the marker from training (one seed) shifts the models toward content: content-only AUROC rises to 0.872 for the box model and to 0.834 for the pixel model (0.681 before), while their render-only response stays lower (0.713, 0.638). Content changes, not the generator's texture, become the main cue once the shortcut is gone.

\begin{table}[t]
\caption{Re-rendering vs.\ content (all validation pairs; standard U-Net rows: mean of three seeds, other rows seed 42): mean probability inside the box and image-level AUROC against authentic images.}
\label{tab:render}
\centering\footnotesize
\setlength{\tabcolsep}{3pt}
\begin{tabular}{@{}lcccccc@{}}
\toprule
 & \multicolumn{3}{c}{mean $p$ in box} & \multicolumn{3}{c}{AUROC vs.\ authentic} \\
\cmidrule(lr){2-4}\cmidrule(lr){5-7}
Model & forged & render & content & forged & render & content \\
\midrule
U-Net \GTb & 0.86 & 0.24 & 0.35 & 0.971 & 0.637 & 0.733 \\
U-Net \GTp & 0.41 & 0.05 & 0.15 & 0.942 & 0.554 & 0.681 \\
DeepLabV3+ \GTb & 0.87 & 0.28 & 0.42 & 0.972 & 0.644 & 0.755 \\
repaired \GTb & 0.83 & 0.36 & 0.59 & 0.947 & 0.713 & 0.872 \\
repaired \GTp$_{\mathrm{nm}}$ & 0.34 & 0.10 & 0.29 & 0.869 & 0.638 & 0.834 \\
\bottomrule
\end{tabular}
\end{table}


\subsection{Held-out generator (H3), now with three seeds}
\label{sec:res-heldout}
Training on Gemini only and testing on the Ideogram edits of validation and test ($n{=}126$ under \GTb, 41 pairs under \GTp) confirms our earlier single-seed result with three seeds (Table~\ref{tab:heldout}). Every Gemini-only model collapses on Ideogram; the \GTb{} model falls from 0.909 to 0.412 per-image AP and the \GTp{} model from 0.567 to 0.264 (both scored under \GTb). The relative drop is smaller for pixel supervision, but its absolute Ideogram score is lower, and under its own ground truth the pixel model is no better on Ideogram than the box model (0.279 vs.\ 0.259). Pixel-change supervision does not, by itself, buy generator transfer. Because Ideogram forgeries almost never carry the marker (green residue in the box ring of 2.1\,\% of Ideogram vs.\ 85.9\,\% of Gemini forgeries), part of this gap may be a marker gap, which this design cannot separate: we did not train Gemini-only models on marker-repaired images.

\begin{table}[t]
\caption{Held-out generator: per-image AP with 95\,\% group-bootstrap CI, Ideogram vs.\ source-matched Gemini (val+test, mean over three seeds).}
\label{tab:heldout}
\centering\footnotesize
\setlength{\tabcolsep}{3pt}
\begin{tabular}{@{}llcc@{}}
\toprule
Training & Scored & Ideogram & Gemini (matched) \\
\midrule
All tools, \GTb & \GTb & 0.629 \ci{0.569}{0.692} & 0.912 \ci{0.896}{0.928} \\
Gemini only, \GTb & \GTb & 0.412 \ci{0.341}{0.481} & 0.909 \ci{0.892}{0.927} \\
Gemini only, \GTp & \GTb & 0.264 \ci{0.213}{0.313} & 0.567 \ci{0.550}{0.583} \\
Gemini only, \GTb & \GTp & 0.259 \ci{0.168}{0.350} & 0.349 \ci{0.322}{0.377} \\
Gemini only, \GTp & \GTp & 0.279 \ci{0.182}{0.382} & 0.577 \ci{0.546}{0.608} \\
Gemini only, \GTg & \GTp & 0.182 \ci{0.121}{0.246} & 0.388 \ci{0.363}{0.413} \\
\bottomrule
\end{tabular}
\end{table}

\subsection{Supervision variants, training length and robustness}
\label{sec:res-variants}
\emph{Joint supervision} (pixel mask plus an auxiliary box channel, seed 42) is the only variant that improves the pixel model on almost every axis we measured (the 89-pair test pixel AUPRC is unchanged, 0.650 vs.\ 0.654): validation AUPRC under \GTp{} rises from 0.610 to 0.622 and under \GTb{} from 0.579 to 0.647, test \GTb{} AUPRC from 0.638 to 0.682, and the blur penalty shrinks from $-59$\,\% to $-43$\,\% (Table~\ref{tab:dualgt}; blur from the robustness runs). Image-level discrimination against authentic images is unchanged (0.897 vs.\ 0.892). With one seed these gains are suggestive, not established.
\emph{Twenty epochs} do not change the picture: \GTb{} 0.926 (vs.\ 0.923), \GTp{} 0.594 (vs.\ 0.610), \GTg{} 0.645 (vs.\ 0.676); the concern that ten epochs were too few is not supported.
\emph{Test-time preprocessing.} Converting the input to grayscale is the most damaging operation we tested, costing 47--67\,\% of the AUPRC of every segmenter for which it was measured (seed-42 models, including the marker-repaired ones) and 59\,\% for the detector, so colour carries evidence beyond the green line. CLAHE changes AUPRC by at most 3\,\%. JPEG Q75, $0.75\times$ rescaling and $\sigma{=}1$ blur hurt pixel and glyph models more than box models (e.g.\ blur $-59$\,\% vs.\ $-21$\,\% for the U-Net), and the marker-repaired models are less robust to JPEG ($-32$ to $-50$\,\%) than their standard counterparts.

\subsection{Box detection versus segmentation}
\label{sec:res-yolo}
A tiled YOLO11s trained on the same boxes is the strongest model in the study \emph{under the box metric} and is no better than the box U-Net under the pixel metric (Tables~\ref{tab:dualgt} and~\ref{tab:yolo}). Over three seeds it reaches 0.965 validation and 0.969 test AUPRC under \GTb{} (U-Net 0.923 / 0.950), test F1@$\tau^*$ 0.946 and IoU 0.918 (U-Net 0.846 / 0.789) and box AP@0.5 0.880 (U-Net 0.873 for seeds 42/43; 0.391 for seed 44, whose calibrated $\tau^*{=}0.95$ fragments the predicted region), while its \GTp{} AUPRC is 0.368 (U-Net \GTb{} model 0.374). The detector depends far less on the marker: its test AP on the 185 residue-free images is 0.878 (U-Net 0.690) and 0.886 on residue-free Gemini images (0.722). It also transfers better across generators: trained on all tools it reaches 0.872 AP on Ideogram (U-Net 0.629; both three seeds), and trained on Gemini only it keeps 0.752 (one seed; U-Net 0.412 over three seeds; Gemini 0.938 vs.\ 0.909). The probes (seed 42 for the detector) suggest why: removing the residue costs the detector only 0.062 AP (U-Net 0.222), whereas the \emph{render-only} input, which keeps the generator's re-rendering but no changed glyph, is still separated from authentic images with AUROC 0.865 (U-Net 0.637). The detector has learned what a regenerated field looks like, a cue that is shared by both generators but says nothing about which pixels changed.

The ablations (seed 42) are only partly as expected. Resizing whole images to 1280 instead of tiling at native resolution costs 0.037 test AUPRC and 0.122 AP on residue-free images, but \emph{raises} box AP@0.5 from 0.88 to 0.96 (0.93 on fields below 1\,\% of the image): a plausible cause is that fields cut at tile borders yield partial boxes that cross-tile NMS does not merge, although the untiled detector also emits about five boxes per image, so the mechanism is not established; tiling helps coverage but not box precision here; model size matters little (n/s/m: 0.963/0.969/0.969 test AUPRC); Ultralytics' default augmentation is as good as the conservative profile (0.971), contrary to our expectation that mosaic and scale jitter would destroy forensic traces; one extra authentic tile per image changes nothing (0.971). The detector is insensitive to blur, CLAHE and rescaling ($\le$1.4\,\% AUPRC) but loses 49\,\% under JPEG Q75 and 59\,\% under grayscale. Contrary to our expectation, detection was not cheaper in our implementation: full-image inference takes 76--98\,ms for the tiled n/s/m detectors, 303\,ms for the 1280 detector, 81\,ms for the U-Net and 122\,ms for DeepLabV3+ (mean 3.1\,MP per image), because per-tile pre- and post-processing dominates.

\begin{table}[t]
\caption{Box detector vs.\ segmenters on the official test set (812 images, \GTb). AUPRC: pooled pixel AUPRC (box painted with its confidence); AP: per-image AP on images without / with marker residue; AP$_{50}$: box AP@0.5 (all images). Seeds: YOLO11s and U-Net 3, others 1. ms: full-image inference on an idle RTX 5080.}
\label{tab:yolo}
\centering\footnotesize
\setlength{\tabcolsep}{2.4pt}
\begin{tabular}{@{}lcccccc@{}}
\toprule
Model & AUPRC & F1@$\tau^*$ & AP\,no-mk & AP\,mk & AP$_{50}$ & ms \\
\midrule
U-Net \GTb{} (3) & 0.950 & 0.846 & 0.690 & 0.972 & 0.873$^\dagger$ & 81 \\
U-Net \GTp{} (3) & 0.638 & 0.616 & 0.503 & 0.639 & 0.107 & 81 \\
DeepLabV3+ \GTb & 0.952 & 0.833 & 0.702 & 0.962 & 0.833 & 122 \\
U-Net \GTb, repaired & 0.939 & 0.797 & 0.701 & 0.920 & 0.805 & 81 \\
\midrule
YOLO11s tiles (3) & \textbf{0.969} & \textbf{0.946} & 0.878 & 0.954 & 0.880 & 98 \\
YOLO11n tiles & 0.963 & 0.926 & 0.812 & 0.946 & 0.871 & 81 \\
YOLO11m tiles & \textbf{0.969} & \textbf{0.947} & \textbf{0.882} & 0.954 & 0.898 & 76 \\
YOLO11s full 1280 & 0.932 & 0.909 & 0.756 & 0.930 & \textbf{0.960} & 303 \\
\midrule
\multicolumn{7}{@{}l}{\scriptsize $^\dagger$seeds 42/43; seed 44 ($\tau^*{=}0.95$) 0.391. Zero-shot TruFor: F1@0.5 0.403.}\\
\bottomrule
\end{tabular}
\end{table}

